% ---------------------------------------------------------------------
\subsection{Discussion: What the Evidence Adds}\label{sec:discussion}
Each guarantee has a well-known ancestor (Section~\ref{sec:related}). For each one we ask what our evidence shows beyond that ancestor, and where it shows nothing.

\emph{Currentness of content, not only of reads (Theorem~\ref{thm:safety}; RQ1, RQ5).} OCC validates at commit that the items a transaction read are unchanged~\cite{occ}. Recent transactional proposals for LLM agents carry this check, or compensation, over to agent actions~\cite{chang2025sagallm,lyu2026coagent,yang2026concurrency}. Such a check corresponds to our request-identity defense over the data a job reads. RQ1 shows that it is not sufficient for generated content. When the contract or the selected context was left out of the key, while all source reads were still checked, stale values became usable in \ContractStalePct\% and \ContextStalePct\% of runs. A generated artifact becomes stale when anything that shaped the question changes: the instruction, the validator, the selected evidence or the owner. RQ5 shows that this matters with real models. Of \LlmStaleTrials{} replies that arrived after their requirement changed, \LlmStaleOldValid{} would have passed the old validator, and \LlmStaleWrongPct\% of those were wrong under the new one. A late reply is therefore not merely outdated; it is usually incorrect. A design that asks the producing model whether a conflicting change invalidates its output~\cite{lyu2026coagent} would delegate this judgment to the component whose replies were wrong.

\emph{Early cutoff under nondeterminism (Theorems~\ref{thm:local} and~\ref{thm:cost}; RQ2, RQ5, RQ6).} Early cutoff is standard in build systems~\cite{mitchell2012shake,mokhov,erdweg2015pluto,hammer2014adapton}, where deterministic tasks make it pay off reliably. With stochastic producers, equality-based cutoff stays sound, because an equal value is the same evidence, but its benefit becomes an empirical question. Neither the build-system literature nor studies of LLM nondeterminism~\cite{ouyang2025nondeterminism} answer it. RQ2 isolates the structural ceiling: with deterministic proposers, cutoff reduced a cosmetic edit of a shared model from \CoupledCosmeticNoCut{} to \CoupledCosmeticExact{} requests. RQ5 and RQ6 measure how much of that ceiling survives real generation. Narrow, canonical slots kept most of it: regenerated designs were value-equal in \LlmCutoffDesignValuePct\% of trials, but raw text was identical in only \LlmCutoffDesignRawPct\%. Free-form code lost most of it. Only \BenchValueEqualPct\% of regenerated benchmark functions were value-equal, although \BenchBehEqualPct\% were behaviorally equal on all plus inputs. The gap between these two numbers is the saving that a semantic cache~\cite{bang2023gptcache,schroeder2026vcache} or a contract-level normal form could recover, at the price of false hits, which exact keys exclude by construction. Slot typing is thus a design parameter, and our data quantify its effect.

\emph{Bounded repair, not bounded retries (Theorems~\ref{thm:single} and~\ref{thm:termination}, Proposition~\ref{prop:eager}).} Feedback-driven retry is a standard ingredient of LLM code generation~\cite{madaan2023selfrefine,shinn2023reflexion,chen2024selfdebug,huang2023agentcoder,schafer2024testpilot}, usually with a per-run iteration cap. Our budget belongs to a dependency episode, which is why the terminal outcome holds under arbitrary interleavings of edits. The upstream gate is why a change starts at most one episode per affected job. The scheduling experiment shows that the gate is not an implementation detail. Without it, the same lattices needed up to \SchedEagerReqRatioMax$\times$ the requests and gained no latency. RQ5 shows the budget at work with real models: retries with validator feedback recovered \LlmFeedbackRecoveredPct\% of failed first attempts. The remaining failures (\LlmJobsFailedPct\% of jobs) ended in explicit verdicts that blocked their dependents, never in silent acceptance or unbounded retry. Section~\ref{sec:rq5} discusses what this means for weaker models, for which most workspaces ended blocked.

\emph{Locally valid is not jointly valid (global predicate, P7; Proposition~\ref{prop:epochs}).} Benchmarks of parallel coding agents report patches that pass alone and fail together~\cite{xia2026stale,khatua2026cooperbench}. This is the LLM counterpart of the result from model consistency that pairwise consistency does not imply global consistency~\cite{stevens}. Our barrier checks the global predicate before publication, and RQ1 shows that removing it publishes inconsistent snapshots (\GlobalPublishPct\% of runs). With real models, however, the barrier never fired in our workflow (\LlmBarrierConflicts{} conflicts in \LlmTrials{} trials): every model chose the same identifier type for all services. Our evidence for the barrier and for the feedback-epoch bound is therefore the adversarial campaign and the model checker, not RQ5. How often teams of real agents produce jointly inconsistent artifacts depends on the workflow, and the studies cited above suggest that richer interfaces make it common.

\emph{Validator-relative correctness (Corollary~\ref{cor:domain}; RQ6).} It is known that weak tests accept wrong programs~\cite{evalplus}, and multi-agent code generators already pair coders with testers~\cite{dong2024selfcollab,huang2023agentcoder,chen2023codet}. RQ6 adds a clean attribution inside a coordinated workflow. Coordination neither created nor hid incorrectness: \BenchAcceptedWrongPct\% of accepted values failed the hidden plus tests, a figure set by the strength of the base tests that served as validator. The corollary states this dependence, and the experiment measures its size.

\emph{Verified once, not per task (Theorems~\ref{thm:refinement} and~\ref{thm:reconfig}; RQ4).} MAS model checkers verify agents~\cite{lomuscio2017mcmas,dennis2012ajpf}, and recent work verifies LLM-agent protocols or workflows per task~\cite{xia2026tracefix,wang2026lean4agent}. Our verification object is the task-independent kernel, and this answers the question of who writes the specifications. A user of the kernel writes jobs, read sets and validators, which are executable checks, as in RQ5 and RQ6, not TLA$^+$. The model and proofs are written once by the kernel's developers, as for verified infrastructure in industry~\cite{newcombe2015aws,howard2025smartcasual}. Trace validation~\cite{cirstea2024trace} is the established path for keeping the implementation aligned with the model. Studies of LLM-written TLA$^+$ report low semantic correctness~\cite{bisharat2026tlallm}, which supports this division of labor rather than asking users or models to specify each workflow.
